\documentclass[oneside,openany,phd]{TMU-Thesis}


% ==================== structure/preamble.tex ====================
% مشخصات پایان‌نامه را در فایلهای faTitle و enTitle وارد نمایید.
% در اینجا برای یکپارچگی، همینجا وارد می‌کنیم

% مشخصات فارسی
\university{تربیت مدرس}
\faculty{مهندسی برق و کامپیوتر}
\department{گروه هوش مصنوعی و رباتیکز}
\subject{مهندسی کامپیوتر}
\field{مهندسی نرم‌افزار}
\credits{---} % در رساله دکتری معمولا به صورت واحدی ذکر نمی‌شود یا فیلد خالی است
\title{یادگیری تقویتی عمیق با شکل‌دهی تصادفی پاداش برای مدیریت عدم قطعیت در سیستم‌های خودوفق}
\firstsupervisor{دکتر ...}
% \secondsupervisor{استاد راهنمای دوم}
\firstadvisor{دکتر ...}
% \secondadvisor{استاد مشاور دوم}
\name{علیرضا}
\surname{کاویانی‌فر}
\studentID{...}
\thesisdate{بهمن ۱۴۰۴}

% مشخصات انگلیسی
\latinuniversity{Tarbiat Modares University}
\latinfaculty{Faculty of Electrical and Computer Engineering}
\latinsubject{Computer Engineering}
\latinfield{Software Engineering}
\latintitle{Randomized Reward-Shaped Deep Reinforcement Learning for Managing Uncertainty in Self-Adaptive Systems}
\firstlatinsupervisor{Dr. ...}
\firstlatinadvisor{Dr. ...}
\latinname{Alireza}
\latinsurname{Kavianifar}
\latinthesisdate{February 2026}
\latinkeywords{Self-Adaptive Systems, Deep Reinforcement Learning, Randomized Reward Shaping, Uncertainty Management, Formal Verification, UPPAAL-SMC}

% فراخوانی بسته‌ها و دستورات استاندارد
% در این فایل، دستورها و تنظیمات مورد نیاز، آورده شده است.
%-------------------------------------------------------------------------------------------------------------------

% در ورژن جدید زی‌پرشین برای تایپ متن‌های ریاضی، این سه بسته، حتماً باید فراخوانی شود
\usepackage{amsthm,amssymb,amsmath}
% بسته‌ای برای تنطیم حاشیه‌های بالا، پایین، چپ و راست صفحه
\usepackage[top=30mm, bottom=25mm, left=25mm, right=35mm]{geometry}
% بسته‌‌ای برای ظاهر شدن شکل‌ها و تصاویر متن
\usepackage{graphicx}
% بسته‌ای برای رسم کادر
\usepackage{framed} 
% بسته‌‌ای برای چاپ شدن خودکار تعداد صفحات در صفحه «معرفی پایان‌نامه»
\usepackage{lastpage}
% بسته‌ و دستوراتی برای ایجاد لینک‌های رنگی با امکان جهش
% \usepackage[pagebackref=false,colorlinks,linkcolor=black,citecolor=black]{hyperref}
% چنانچه قصد پرینت گرفتن نوشته خود را دارید، خط بالا را غیرفعال و  از دستور زیر استفاده کنید چون در صورت استفاده از دستور زیر‌‌، 
% لینک‌ها به رنگ سیاه ظاهر خواهند شد که برای پرینت گرفتن، مناسب‌تر است
% \usepackage[pagebackref=false]{hyperref}
% بسته‌ لازم برای تنظیم سربرگ‌ها
\usepackage{fancyhdr}
\setlength{\headheight}{16pt}
%
% تنظیمات کنترل جای‌گذاری شناورها (تصاویر و جداول) جهت جلوگیری از رانش به انتهای فصل
\renewcommand{\topfraction}{0.85}
\renewcommand{\bottomfraction}{0.80}
\renewcommand{\textfraction}{0.15}
\renewcommand{\floatpagefraction}{0.75}
\setcounter{topnumber}{3}
\setcounter{bottomnumber}{2}
\setcounter{totalnumber}{4}
%
\usepackage{setspace}
\usepackage{algorithm}
\usepackage{algorithmic}
\usepackage{subfigure}
\usepackage[subfigure]{tocloft}
\usepackage{multirow}
\usepackage{xspace}


% بسته‌ای برای ظاهر شدن «مراجع» و «نمایه» در فهرست مطالب
\usepackage[nottoc]{tocbibind}
% دستورات مربوط به ایجاد نمایه
\usepackage{makeidx}
\makeindex
% تعریف شیم‌های مورد نیاز برای نسخه‌های جدید بسته bidi در برخی محیط‌ها
% Removal of redundant and potentially broken definitions
\makeatletter
% شیم‌های سازگاری برای نسخه‌های جدید لاتک و بسته‌های bidi/xepersian
\providecommand{\IfClassLoadedT}[2]{\@ifclassloaded{#1}{#2}{}}
\providecommand{\IfClassLoadedF}[2]{\@ifclassloaded{#1}{}{#2}}
\providecommand{\IfClassLoadedTF}[3]{\@ifclassloaded{#1}{#2}{#3}}
\providecommand{\IfPackageLoadedT}[2]{\@ifpackageloaded{#1}{#2}{}}
\providecommand{\IfPackageLoadedF}[2]{\@ifpackageloaded{#1}{}{#2}}
\providecommand{\IfPackageLoadedTF}[3]{\@ifpackageloaded{#1}{#2}{#3}}
\providecommand{\@iffileloaded}[3]{\expandafter\ifx\csname ver@#1\endcsname\relax#3\else#2\fi}
\providecommand{\IfFileLoadedT}[2]{\@iffileloaded{#1}{#2}{}}
\providecommand{\IfFileLoadedF}[2]{\@iffileloaded{#1}{}{#2}}
\providecommand{\IfFileLoadedTF}[3]{\@iffileloaded{#1}{#2}{#3}}
\makeatother

% Standalone configuration to support individual chapter compilation
\ifdefined\ifstandalone
\else
  \newif\ifstandalone
  \standalonetrue
\fi

% Cross-document label resolution for standalone chapter builds
\ifstandalone
  \usepackage{xr-hyper}
\fi

%%%%%%%%%%%%%%%%%%%%%%%%%%
% فراخوانی بسته زی‌پرشین و تعریف قلم فارسی و انگلیسی
\usepackage[pagebackref=false,colorlinks,linkcolor=black,citecolor=black]{hyperref}

\ifstandalone
  \IfFileExists{master_labels.aux}{%
    \externaldocument{master_labels}%
  }{}
\fi
\usepackage{notoccite}

\ExplSyntaxOn
\cs_if_exist:NTF \g__hyp_linknestlevel_int { } { \int_new:N \g__hyp_linknestlevel_int }
\cs_if_exist:NTF \l__hyp_annot_GoTo_bool { } { \bool_new:N \l__hyp_annot_GoTo_bool }
\cs_if_exist:NTF \l__hyp_annot_URI_bool { } { \bool_new:N \l__hyp_annot_URI_bool }
\cs_if_exist:NTF \l__hyp_href_url_ismap_bool { } { \bool_new:N \l__hyp_href_url_ismap_bool }
\cs_if_exist:NTF \l__hyp_text_enc_uri_print_tl { } { \tl_new:N \l__hyp_text_enc_uri_print_tl }
\cs_if_exist:NTF \l__hyp_uri_tmpa_tl { } { \tl_new:N \l__hyp_uri_tmpa_tl }

% Defining missing l3pdf/pdfdict commands if not present
\cs_if_exist:NTF \pdfdict_put:nno { } { \cs_new:Npn \pdfdict_put:nno #1#2#3 { } }
\cs_if_exist:NTF \pdfdict_put:nnn { } { \cs_new:Npn \pdfdict_put:nnn #1#2#3 { } }
\cs_if_exist:NTF \pdfdict_use:n { } { \cs_new:Npn \pdfdict_use:n #1 { } }
\cs_if_exist:NTF \pdfannot_dict_put:nne { } { \cs_new:Npn \pdfannot_dict_put:nne #1#2#3 { } }
\cs_if_exist:NTF \pdfannot_link:nen { } { \cs_new:Npn \pdfannot_link:nen #1#2#3 { } }
\cs_if_exist:NTF \__hyp_check_link_nesting:TF { } { \cs_new:Npn \__hyp_check_link_nesting:TF #1#2 { #1 } }
\cs_if_exist:NTF \__hyp_text_pdfstring:eoN { } { \cs_new:Npn \__hyp_text_pdfstring:eoN #1#2#3 { } }
\ExplSyntaxOff

% تعریف دستورات گمشده برای جلوگیری از خطای renewcommand در بسته‌های داخلی
\makeatletter
\providecommand{\foottextfont}{}
\providecommand{\LTRfoottextfont}{}
\providecommand{\RTLfoottextfont}{}
\providecommand{\abstractname}{}
\providecommand{\latinabstract}{}
\providecommand{\AutoPersianDigits}{}
\makeatother

\usepackage{xepersian}
\settextfont[Path=font/,Extension=.ttf,Scale=1,
    BoldFont=XB NiloofarBd,
    ItalicFont=XB NiloofarIt,
    BoldItalicFont=XB NiloofarBdIt
]{XB Niloofar}
\setlatintextfont[Scale=0.9]{Times New Roman}

%%%%%%%%%%%%%%%%%%%%%%%%%%
% چنانچه می‌خواهید اعداد در فرمول‌ها، انگلیسی باشد، خط زیر را غیرفعال کنید
\setdigitfont[Path=font/,Extension=.ttf,Scale=1,
    BoldFont=XB NiloofarBd,
    ItalicFont=XB NiloofarIt,
    BoldItalicFont=XB NiloofarBdIt
]{XB Niloofar}
%%%%%%%%%%%%%%%%%%%%%%%%%%
% تعریف قلم‌های فارسی و انگلیسی اضافی برای استفاده در بعضی از قسمت‌های متن
\defpersianfont\titlefont[Path=font/,Extension=.ttf,Scale=1]{XB Titre}
% \defpersianfont\iranic[Scale=1.1]{XB Zar Oblique}%Italic}%
% \defpersianfont\nastaliq[Scale=1.2]{IranNastaliq}

%%%%%%%%%%%%%%%%%%%%%%%%%%
% دستوری برای حذف کلمه «چکیده»
\renewcommand{\abstractname}{}
% دستوری برای حذف کلمه «abstract»
%\renewcommand{\latinabstract}{}
% دستوری برای تغییر نام کلمه «اثبات» به «برهان»
\renewcommand\proofname{\textbf{برهان}}
% دستوری برای تغییر نام کلمه «کتاب‌نامه» به «مراجع»
\AtBeginDocument{%
  \renewcommand{\bibname}{مراجع}%
  \ifx\settocbibname\undefined\else\settocbibname{مراجع}\fi
}
% دستوری برای تعریف واژه‌نامه انگلیسی به فارسی
\newcommand\persiangloss[2]{#1\dotfill\lr{#2}\\}
% دستوری برای تعریف واژه‌نامه فارسی به انگلیسی 
\newcommand\englishgloss[2]{#2\dotfill\lr{#1}\\}
% تعریف دستور جدید «\پ» برای خلاصه‌نویسی جهت نوشتن عبارت «پروژه/پایان‌نامه/رساله»
\newcommand{\پ}{پروژه/پایان‌نامه/رساله }

%\newcommand\BackSlash{\char`\\}

%%%%%%%%%%%%%%%%%%%%%%%%%%
\SepMark{-}

% تعریف و نحوه ظاهر شدن عنوان قضیه‌ها، تعریف‌ها، مثال‌ها و ...
\theoremstyle{definition}
\newtheorem{definition}{تعریف}[section]
\theoremstyle{plain}
\newtheorem{theorem}[definition]{قضیه}
\newtheorem{lemma}[definition]{لم}
\newtheorem{proposition}[definition]{گزاره}
\newtheorem{corollary}[definition]{نتیجه}
\newtheorem{remark}[definition]{ملاحظه}
\theoremstyle{definition}
\newtheorem{example}[definition]{مثال}

%\renewcommand{\theequation}{\thechapter-\arabic{equation}}
%\def\bibname{مراجع}
\numberwithin{algorithm}{chapter}
\def\listalgorithmname{فهرست الگوریتم‌ها}
\def\listfigurename{فهرست تصاویر}
\def\listtablename{فهرست جداول}

%%%%%%%%%%%%%%%%%%%%%%%%%%%%
% دستورهایی برای سفارشی کردن سربرگ صفحات
% \newcommand{\SetHeader}{
% \csname@twosidetrue\endcsname
% \pagestyle{fancy}
% \fancyhf{} 
% \fancyhead[OL,EL]{\thepage}
% \fancyhead[OR]{\small\rightmark}
% \fancyhead[ER]{\small\leftmark}
% \renewcommand{\chaptermark}[1]{%
% \markboth{\thechapter-\ #1}{}}
% }
%%%%%%%%%%%%5
%\def\MATtextbaseline{1.5}
%\renewcommand{\baselinestretch}{\MATtextbaseline}
\doublespacing
%%%%%%%%%%%%%%%%%%%%%%%%%%%%%
% دستوراتی برای اضافه کردن کلمه «فصل» و شماره ترتیبی حروفی در فهرست مطالب

\makeatletter
% سفارشی‌سازی سیستم ابجد تک‌حرفی (\Abjad) برای نمایش «الف» به جای «آ»
\def\@Abjad#1{%
  \ifcase #1\or الف\or ب\or ج\or د\or ه\or و\or ز\or ح\or ط%
  \or ی\or ک\or ل\or م\or ن\or س\or ع\or ف\or ص\or ق\or ر\or ش\or ت\or ث\or خ%
  \or ذ\or ض\or ظ\or غ\else \@ctrerr \fi}

% سفارشی‌سازی عدد ۱ در سیستم ابجد برای نمایش کلمه «الف» به جای «آ»
\def\@abjad#1{%
\ifnum#1=1 الف%
\else
  \ifnum#1>1999 \xpg@warning{Illegal value (#1) for abjad numeral} {#1}%
  \else
    \ifnum#1<\z@\space\xpg@warning{Illegal value (#1) for abjad numeral}%
    \else
      \ifnum#1<10\expandafter\abj@num@i\number#1%
      \else
        \ifnum#1<100\expandafter\abj@num@ii\number#1%
        \else
          \ifnum#1<\@m\expandafter\abj@num@iii\number#1%
          \else
            \ifnum#1<\@M\expandafter\abj@num@iv\number#1%
            \fi
          \fi
        \fi
      \fi
    \fi
  \fi
\fi}
\def\abj@num@i#1{%
  \ifcase#1\or ا\or ب\or ج\or د%
           \or ه\or و\or ز\or ح\or ط\fi
  \ifnum#1=\z@\abjad@zero\fi}

\newcommand{\persianchapnum}[1]{%
  \edef\@arg{#1}%
  \def\@one{1}\def\@pone{۱}%
  \def\@two{2}\def\@ptwo{۲}%
  \def\@three{3}\def\@pthree{۳}%
  \def\@four{4}\def\@pfour{۴}%
  \def\@five{5}\def\@pfive{۵}%
  \def\@six{6}\def\@psix{۶}%
  \def\@seven{7}\def\@pseven{۷}%
  \def\@eight{8}\def\@peight{۸}%
  \def\@nine{9}\def\@pnine{۹}%
  \def\@ten{10}\def\@pten{۱۰}%
  \ifx\@arg\@one اول%
  \else\ifx\@arg\@pone اول%
  \else\ifx\@arg\@two دوم%
  \else\ifx\@arg\@ptwo دوم%
  \else\ifx\@arg\@three سوم%
  \else\ifx\@arg\@pthree سوم%
  \else\ifx\@arg\@four چهارم%
  \else\ifx\@arg\@pfour چهارم%
  \else\ifx\@arg\@five پنجم%
  \else\ifx\@arg\@pfive پنجم%
  \else\ifx\@arg\@six ششم%
  \else\ifx\@arg\@psix ششم%
  \else\ifx\@arg\@seven هفتم%
  \else\ifx\@arg\@pseven هفتم%
  \else\ifx\@arg\@eight هشتم%
  \else\ifx\@arg\@peight هشتم%
  \else\ifx\@arg\@nine نهم%
  \else\ifx\@arg\@pnine نهم%
  \else\ifx\@arg\@ten دهم%
  \else\ifx\@arg\@pten دهم%
  \else #1\fi\fi\fi\fi\fi\fi\fi\fi\fi\fi\fi\fi\fi\fi\fi\fi\fi\fi\fi\fi
}

\newlength\mylenprt
\newlength\mylenchp
\newlength\mylenapp

\renewcommand\cftpartpresnum{\partname~}
\renewcommand\cftchappresnum{\chaptername~}
\renewcommand\cftchapaftersnum{:}

\settowidth\mylenprt{\cftpartfont\cftpartpresnum\cftpartaftersnum}
% تنظیم پهنای کادر شماره متناسب با عریض‌ترین عنوان فصل («فصل چهارم: »)
\settowidth\mylenchp{\cftchapfont\chaptername~چهارم:\space}
\settowidth\mylenapp{\cftchapfont\appendixname~\cftchapaftersnum}
\addtolength\mylenprt{\cftpartnumwidth}
\addtolength\mylenchp{0.5em}
\addtolength\mylenapp{\cftchapnumwidth}

\setlength\cftpartnumwidth{\mylenprt}
\setlength\cftchapnumwidth{\mylenchp}

% اعمال تبدیل شماره فصل به حروف در فهرست مطالب با رعایت چیدمان راست‌به‌چپ (RTL)
\renewcommand*{\l@chapter}[2]{%
  \ifnum \c@tocdepth >\m@ne
    \addpenalty{-\@highpenalty}%
    \vskip \cftbeforechapskip
    {\if@RTL\rightskip\else\leftskip\fi \cftchapindent\relax
     \if@RTL\leftskip\else\rightskip\fi \@tocrmarg
     \parfillskip -\if@RTL\leftskip\else\rightskip\fi
     \parindent \cftchapindent\relax\@afterindenttrue
     \interlinepenalty\@M
     \leavevmode
     \@tempdima \cftchapnumwidth\relax
     \let\@cftbsnum \cftchappresnum
     \let\@cftasnum \cftchapaftersnum
     \let\@cftasnumb \cftchapaftersnumb
     \advance\if@RTL\rightskip\else\leftskip\fi \@tempdima \null\nobreak\hskip -\if@RTL\rightskip\else\leftskip\fi
     {\cftchapfont \def\numberline##1{\hb@xt@\@tempdima{\@cftbsnum \persianchapnum{##1}\@cftasnum\hfil}\@cftasnumb}#1}\nobreak
     \cftchapfillnum{#2}}%
  \fi}

% فعال‌سازی نقطه‌چین برای فصول در فهرست مطالب هماهنگ با بخش‌ها
\renewcommand{\cftchapleader}{\cftdotfill{\cftsecdotsep}}

% اصلاح نمایش پیوست‌ها در فهرست مطالب (تبدیل «فصل آ» به «پیوست آ»)
\newif\ifinappendix \inappendixfalse
\g@addto@macro\appendix{%
  \inappendixtrue
  \addtocontents{toc}{%
    \protect\renewcommand{\protect\cftchappresnum}{\protect\appendixname~}%
  }%
}

% اصلاح سربرگ بالای صفحات (Header) با نام ترتیبی حروفی فصول
\renewcommand{\chaptermark}[1]{%
  \ifinappendix
    \markboth{\appendixname\space\thechapter:\space#1}{}%
  \else
    \begingroup
      \ifcase\c@chapter
        \def\@tempchap{}%
        \or\def\@tempchap{اول}%
        \or\def\@tempchap{دوم}%
        \or\def\@tempchap{سوم}%
        \or\def\@tempchap{چهارم}%
        \or\def\@tempchap{پنجم}%
        \or\def\@tempchap{ششم}%
        \or\def\@tempchap{هفتم}%
        \or\def\@tempchap{هشتم}%
        \or\def\@tempchap{نهم}%
        \or\def\@tempchap{دهم}%
        \else\def\@tempchap{\thechapter}%
      \fi
      \expandafter\endgroup
      \expandafter\markboth\expandafter{\expandafter\chaptername\expandafter\space\@tempchap:\space#1}{}%
  \fi
}

% اصلاح سرتیتر شروع فصول در متن اصلی به صورت حروف فارسی
\def\@makechapterhead#1{%
  \vspace*{50\p@}%
  {\parindent \z@ \raggedright \normalfont
    \ifnum \c@secnumdepth >\m@ne
      \if@mainmatter
        \huge\bfseries \@chapapp\space 
        \ifx\@chapapp\appendixname
          \thechapter
        \else
          \persianchapnum{\the\c@chapter}
        \fi
        \par\nobreak
        \vskip 20\p@
      \fi
    \fi
    \interlinepenalty\@M
    \Huge \bfseries #1\par\nobreak
    \vskip 40\p@
  }}
\makeatother

\makeatletter
{\def\thebibliography#1{\chapter*{\refname\@mkboth
   {\uppercase{\refname}}{\uppercase{\refname}}}\list
   {[\arabic{enumi}]}{\settowidth\labelwidth{[#1]}
   \rightmargin\labelwidth
   \advance\rightmargin\labelsep
   \advance\rightmargin\bibindent
   \itemindent -\bibindent
   \listparindent \itemindent
   \parsep \z@
   \usecounter{enumi}}
   \def\newblock{}
   \sloppy
   \sfcode`\.=1000\relax}}
\makeatother

% Robust macros for RS-DRL+ and architectural patterns
\xspaceaddexceptions{، ؛ » \footnote \LTRfootnote ^^^^200c}
\newcommand{\rsdrlplus}{\lr{RS-DRL\textsuperscript{+}}\xspace}
\newcommand{\sgos}{\lr{SG-OS}\xspace}
\newcommand{\caepv}{\lr{CA-EPV}\xspace}
\newcommand{\pgalr}{\lr{PG-ALR}\xspace}
\newcommand{\mursdrl}{\lr{MU-RS-DRL}\xspace}
\newcommand{\cmrsdrl}{\lr{CM-RS-DRL}\xspace}


% تنظیمات اختصاصی برای فهرست مطالب (TOC)
\setcounter{secnumdepth}{3}
\setcounter{tocdepth}{3}

\raggedbottom

% تنظیمات نقطه‌چین برای فهرست مطالب
\usepackage{tocloft}
\providecommand{\cftchapdotfill}{\cftdotfill{\cftdotsep}}
\providecommand{\cftsecdotfill}{\cftdotfill{\cftdotsep}}
\providecommand{\cftsubsecdotfill}{\cftdotfill{\cftdotsep}}


% تنظیمات سربرگ برای رفع مشکل نمایش حروف فارسی (جعبه‌های خالی)
\usepackage{fancyhdr}
\pagestyle{fancy}
\fancyhf{}
\fancyhead[L]{\thepage}
\fancyhead[R]{\small\leftmark}
\renewcommand{\chaptermark}[1]{\markboth{\chaptername\ \thechapter.\ #1}{}}
\renewcommand{\sectionmark}[1]{\markright{\thesection.\ #1}}
% حذف حالت ایتالیک/کج که باعث به هم ریختگی فونت فارسی در سربرگ می‌شود
\let\oldheadrule\headrule
\renewcommand{\headrule}{\oldheadrule}
\fancyhead[R]{\small\nouppercase{\leftmark}}



\begin{document}
\AutoPersianDigits
\setstretch{1.5}

\appendix

% =========================================================================
% پیوست الف: جریان‌های کاری کامل فرآیند
% =========================================================================
\chapter{جریان‌های کاری کامل فرآیند}
\label{chap:appendix_workflows}

این پیوست به تشریح جزئیات جریان‌های داده، تعاملات میان‌مؤلفه‌ای و نمودارهای فعالیت فرآیند انطباق در سطوح مختلف معماری \lr{RS-DRL} و معماری دوسطحی \rsdrlplus اختصاص دارد. هدف از ارائه این نمودارها و فرآیندهای گام‌به‌گام، مستندسازی دقیق جریان اجرای نرم‌افزاری جهت بازتولیدپذیری مهندسی در محیط‌های عملیاتی است.

\section{معماری یکپارچه و تفکیک زمانی در چرخه \texorpdfstring{\lr{MAPE-K}}{MAPE-K}}
\label{sec:app_mape_k_integration}

در حلقه‌های خودوفق سنتی، اعتبارسنجی صوری و تصمیم‌گیری به صورت متوالی و مسدودکننده در مسیر مستقیم کنترل اجرا می‌شوند که به نقض مهلت‌های زمانی در سیستم‌های بلادرنگ می‌انجامد. در معماری پیشنهادی \lr{RS-DRL}، تفکیک زمانی میان مسیر تصمیم‌گیری مستقیم برخط و مسیر اعتبارسنجی صوری ناهمگام برقرار گردیده است. شکل \ref{fig:app_mape_k_diagram} جریان تعامل میان مؤلفه‌های چهارگانه حلقه \lr{MAPE-K} و مخزن دانش مرکزی را به تصویر می‌کشد.

\begin{figure}[htbp]
  \centering
  \includegraphics[width=0.92\textwidth]{figures/process-workflow.pdf}
  \caption{جریان داده و کنترل شش‌خط‌شناوری در حلقه \lr{MAPE-K} با تفکیک زمانی میان استنتاج برخط و اعتبارسنجی صوری و یادگیری ناهمگام در مخزن دانش مشترک}
  \label{fig:app_mape_k_diagram}
\end{figure}

همان‌گونه که در شکل \ref{fig:app_mape_k_diagram} مشاهده می‌شود، جریان داده شامل چهار مرحله کلیدی زیر است:

\paragraph{مرحله پایش و انتقال مشاهدات:} کاوشگرهای حسگری مستقر در سیستم هدف، متغیرهای عملیاتی را گردآوری نموده و بردار حالت خام را به مؤلفه پایش ارسال می‌کنند. مؤلفه پایش با اعمال فیلترهای آماری و حذف نویزهای بسامدبالا، بردار حالت نرمال‌شده $s_t$ را تولید و در مخزن دانش ذخیره می‌نماید.

\paragraph{مرحله تحلیل و ارزیابی شرایط تنش:} مؤلفه تحلیل وضعیت جاری را با مشخصه‌های صوری و آستانه‌های بحرانی مقایسه می‌کند. در صورت نقض قید جانشین یا فراتر رفتن شاخص تنش از دروازه مجاز (\lr{$\mathcal{E}(o_f) > \epsilon_{\text{threshold}}$})، سیگنال هشدار فعال شده و درخواست تعیین پیکربندی انطباقی به مؤلفه برنامه‌ریزی مخابره می‌گردد.

\paragraph{مرحله برنامه‌ریزی دوسطحی:} مؤلفه برنامه‌ریزی بر اساس راهبرد دوسطحی عمل می‌کند. در شرایط عادی، پیشنهاد پایه سیاست $a_0$ در زمانی زیر نیم میلی‌ثانیه تولید و بلافاصله به مؤلفه اجرا فرستاده می‌شود. در شرایط تنش محیطی، پالایش محلی هرزمانی بر روی زیرفضای حساسیت محدود فعال می‌گردد.

\paragraph{مرحله اجرا و اعمال فرمان‌های انطباق:} مؤلفه اجرا طرح انطباق انتخاب‌شده را به دستورات پیکربندی سطح سخت‌افزار ترجمه کرده و از طریق واسط‌های راه‌انداز سخت‌افزاری به محرک‌های سیستم اعمال می‌نماید.

\section{سازوکار پایش و ساخت بردار حالت}
\label{sec:app_monitoring_state}

ساخت بردار حالت با ابعاد مشخص و بدون نویز، پیش‌شرط قطعی برای همگرایی شبکه‌های عصبی عمیق در زمان اجرا است. شکل \ref{fig:app_system_context} مرزبندی میان سیستم کنترل خودوفق و سیستم تحت مدیریت را نمایش می‌دهد.

\begin{figure}[htbp]
  \centering
  \includegraphics[width=0.85\textwidth]{figures/system-context.pdf}
  \caption{مرزهای تعامل میان سیستم تحت مدیریت و کنترل‌کننده خودوفق مبتنی بر یادگیری تقویتی عمیق}
  \label{fig:app_system_context}
\end{figure}

فرآیند ساخت بردار حالت در مؤلفه پایش به سه گام پیاپی تفکیک می‌گردد:

\paragraph{گام اول: نمونه‌برداری و همگام‌سازی:}
داده‌های دریافتی از حسگرهای ناهمگن با بسامدهای نمونه‌برداری متفاوت همگام‌سازی شده و وقفه‌های زمانی داده‌ها با روش درون‌یابی خطی پوشش داده می‌شوند.

\paragraph{گام دوم: استخراج ویژگی‌های پویا:}
ویژگی‌های پویا نظیر نرخ تغییرات توان سیگنال، شتاب تغییرات خطا و حاشیه فاصله تا مرزهای بحرانی محاسبه شده و به بردار وضعیت افزوده می‌شوند.

\paragraph{گام سوم: نرمال‌سازی بازه مقادیر بردار حالت:}
تمامی مؤلفه‌های بردار حالت با استفاده از توابع مقیاس‌بندی خطی به بازه $[0, 1]$ نگاشت می‌شوند تا از اشباع گرادیان در لایه‌های ابتدایی شبکه عصبی عمیق پیشگیری به عمل آید.

\section{جریان کاری مرحله برنامه‌ریزی و استنتاج دوسطحی}
\label{sec:app_planning_workflow}

شکل \ref{fig:app_planner_internal} معماری داخلی برنامه‌ریز و مسیرهای سه‌گانه کنترلی را به تفصیل نمایش می‌دهد.

\begin{figure}[htbp]
  \centering
  \includegraphics[width=0.82\textwidth]{figures/fig_decision_hierarchy_flowchart.pdf}
  \caption[معماری سلسله‌مراتبی درونی برنامه‌ریز دوسطحی (\pgalr)]{معماری سلسله‌مراتبی درونی برنامه‌ریز دوسطحی (\pgalr): تفکیک مسیر سریع محاسباتی (\lr{Fast Path})، دروازه تنش و مسیر پالایش محلی هرزمانی}
  \label{fig:app_planner_internal}
\end{figure}

جریان تصمیم‌گیری در برنامه‌ریز دوسطحی از سازوکار کنترلی زیر تبعیت می‌نماید:

\paragraph{مسیر سریع (\lr{Fast Path}):} بردار حالت $s$ مستقیماً به شبکه عصبی عمیق $Q$ (\lr{DQN}) تغذیه شده و اقدام بهینه اسمی $a_0 = \arg\max_{a} Q(s, a)$ با میانه تأخیر $0.485\,\text{ms}$ محاسبه می‌گردد.

\paragraph{دروازه ارزیابی تنش (\lr{Stress Gate}):} اقدام اولیه $a_0$ توسط مدل‌های جانشین سبک‌وزن ارزیابی می‌شود. چنانچه تمامی قیود ارضا شوند، اقدام $a_0$ بدون ورود به فاز پالایش صادر می‌گردد.

\paragraph{مسیر پالایش محلی هرزمانی (\lr{Refinement Path}):} در صورت نقض قید، سازوکار \pgalr فضای جستجو را به زیرفضای حساسیت $m$ بُعدی محدود ساخته و کاندیداهای همسایگی $|\mathcal{N}(a_0)| \le 2m + 1$ را در بازه بودجه زمانی باقیمانده پالایش می‌کند.

\section{سناریوهای مرجع تشریحی چرخه انطباق}
\label{sec:appendix_running_example}

جهت تبیین گام‌به‌گام و ملموس کارکردهای پایش، تحلیل، برنامه‌ریزی و اجرا در سیستم‌های خودوفق، این بخش دو سناریوی مرجع مبتنی بر بستر شبکه‌های حسگر بی‌سیم (\lr{DeltaIoT}) را با جزئیات کامل و ویژگی‌های محاسباتی متفاوت تشریح می‌نماید.

\subsection{سناریوی اول: شبکه دو گره‌ای با قید کیفیت خدمات تک‌هدفه}
\label{subsec:app_scenario_2node}
در این سناریو، چرخه کامل تصمیم‌گیری برخط در قالب یک شبکه حسگر بی‌سیم ساده تبیین می‌شود. شبکه‌ای شامل ۲ گره حسگر ($M_1$ و $M_2$) را در نظر بگیرید که داده‌ها را به یک دروازه مرکزی ارسال می‌کنند و هر گره دارای ۲ گزینه برای تنظیم توان ارسال است: توان کم ($a_{\text{low}}$) و توان زیاد ($a_{\text{high}}$) که مجموعاً $2^2 = 4$ گزینه وفقی را تشکیل می‌دهند. هدف کیفی سیستم این است که نرخ اتلاف بسته‌ها ($P$) همواره کمتر از $10\%$ حفظ شود.

چرخه انطباق با شروع ناگهانی بارندگی یا تداخل محیطی آغاز می‌شود که کیفیت پیوندهای رادیویی را تضعیف کرده و مؤلفه \textbf{پایش}، نرخ اتلاف بسته‌های شبکه را $12\%$ گزارش می‌کند و بردار وضعیت جاری به عنوان $s_t = \langle E, P=0.12, L \rangle$ ساخته می‌شود. در فاز \textbf{تحلیل}، مقدار اتلاف بسته ($0.12$) با آستانه مجاز ($0.10$) مقایسه شده و به دلیل نقض قید ($0.12 > 0.10$)، بلافاصله سیگنال ضرورت انطباق برای فعال‌سازی برنامه‌ریز صادر می‌گردد. در فاز \textbf{برنامه‌ریزی}، بردار حالت $s_t$ به عنوان ورودی به شبکه \lr{DQN} آموزش‌دیده ارائه می‌شود؛ شبکه در یک گذر رو به جلوی مستقیم با دیرند $0.2\,\text{ms}$، مقادیر $Q$ را برای ۴ گزینه موجود محاسبه کرده و گزینه $a^* = \langle a_{\text{high}}, a_{\text{low}} \rangle$ (افزایش توان گره اول و تثبیت توان گره دوم) به عنوان بیشترین ارزش پیش‌بینی‌شده برگزیده می‌شود. نهایتاً در فاز \textbf{اعتبارسنجی صوری و اجرا}، اقدام پیشنهادی $a^*$ به موتور \lr{UPPAAL-SMC} ارسال شده و با تایید ارضای قید احتمالاتی ($r_{\text{smc}} = 1$)، فرمان بازپیکربندی به راه‌اندازهای سخت‌افزاری صادر می‌گردد و نرخ اتلاف بسته به $4\%$ کاهش می‌یابد.

\subsection{سناریوی دوم: شبکه چهار گره‌ای با اهداف متضاد کیفیت خدمات}
\label{subsec:app_scenario_4node}
جهت تبیین عملکرد سیستم در شرایطی که اهداف کیفی در تناقض مستقیم با یکدیگر قرار دارند، یک سناریوی جامع‌تر بررسی می‌شود. شبکه‌ای شامل ۴ گره حسگر ($M_1$ تا $M_4$) در توپولوژی خطی داده‌ها را به دروازه ارسال می‌کنند که هر گره دارای ۳ سطح توان گسسته (پایین، متوسط و بالا) است و در نتیجه فضای وفقی شامل $3^4 = 81$ پیکربندی متفاوت می‌باشد. اهداف کیفی سه قید متضاد را شامل می‌شوند: نرخ اتلاف بسته‌ها باید کمتر از $10\%$ باشد ($P_{\text{net}} < 0.10$)، میانگین تأخیر ارسال باید کمتر از $200\,\text{ms}$ باشد ($L_{\text{net}} < 200\,\text{ms}$)، و مصرف کل انرژی شبکه ($E$) باید کمینه گردد.

در پی بروز اختلال شدید در کانال‌های ارتباطی که کیفیت پیوند گره‌های $M_1$ و $M_2$ را تضعیف می‌کند، حسگرها بردار وضعیت ۱۲ بعدی خام زیر را (با انرژی نرمال‌شده چرخه، نرخ اتلاف کسری، و تأخیر بر حسب میلی‌ثانیه پیش از نرمال‌سازی به $[0, 1]$) پایش و گزارش می‌دهند:
\begin{equation}
\label{eq:app_state_vector_12d}
\begin{aligned}
s_t = \big\langle &E = [0.3,\, 0.3,\, 0.25,\, 0.25], \\
&P = [0.18,\, 0.15,\, 0.04,\, 0.05], \\
&L = [120,\, 110,\, 80,\, 75] \big\rangle
\end{aligned}
\end{equation}
در این وضعیت، با فرض تجمیع وزنی بر پایه سهم بار ترافیکی ارسال‌شده و رله‌شده گره‌ها به سوی دروازه (به‌گونه‌ای که دو گره بالادستی $M_1$ و $M_2$ ترافیک گره‌های پایینی را نیز رله می‌کنند و میانگین وزنی کل از رابطه $\bar{X}_{\text{net}} = \sum_i \omega_i X_i$ حاصل می‌شود)، نرخ اتلاف دو گره اول ($18\%$ و $15\%$) بر عملکرد کل شبکه غلبه یافته و با وزن ترافیکی بالاتر گره‌های رله‌کننده ($M_1, M_2$)، اتلاف مؤثر کل به $P_{\text{net}} = 16.5\%$ افزایش یافته و قید اول را نقض می‌کند (در حالی که میانگین وزنی تأخیر تجمعی در این مسیرها حدود $100$ تا $110\,\text{ms}$ باقی می‌ماند). با ارائه بردار $s_t$ به استنتاج‌گر \lr{DQN}، عامل به ارزیابی مقایسه‌ای گزینه‌های کاندیدا می‌پردازد: اقدام کم‌مصرف ($a_{\text{low}} = \langle 0, 0, 0, 0 \rangle$) که تمام گره‌ها را در توان کم نگه می‌دارد، ارزش نازل $Q = 0.21$ را کسب می‌کند زیرا اتلاف بسته‌ها مهار نمی‌شود؛ در مقابل، اقدام تهاجمی ($a_{\text{brute}} = \langle 2, 2, 2, 2 \rangle$) با افزایش توان همه گره‌ها اتلاف را به $3\%$ می‌رساند، اما به علت تداخل شدید کانال، تأخیر را به $240\,\text{ms}$ رسانده و قید دوم را نقض می‌کند ($Q = 0.42$)؛ در نهایت، اقدام متوازن هوشمند ($a_{\text{balanced}} = \langle 2, 2, 0, 0 \rangle$) با افزایش هدفمند توان گره‌های کلیدی $M_1$ و $M_2$ و حفظ توان پایین سایر گره‌ها، بالاترین ارزش ($Q = 0.92$) را کسب می‌کند، زیرا هم قید اتلاف ($5\%$) و هم قید تأخیر ($155\,\text{ms}$) را با حداقل مصرف انرژی ارضا می‌سازد. با تأیید صوری \lr{UPPAAL-SMC} ($r_{\text{smc}}=1$)، فرمان نهایی به گره‌ها اعمال شده و موازنه بهینه پارتو بازیابی می‌گردد.

\newpage

% =========================================================================
% پیوست ب: لیست کامل الگوریتم‌ها
% =========================================================================
\chapter{لیست کامل الگوریتم‌ها}
\label{chap:appendix_algorithms}

این پیوست متن صوری و شبه‌کد یکپارچه الگوریتم‌های پیشنهادی رساله را برای بازتولید مستقیم در زبان‌های برنامه‌نویسی استاندارد ارائه می‌نماید.

\section{الگوریتم جامع آموزش چارچوب پایه‌ای \texorpdfstring{\lr{RS-DRL}}{RS-DRL}}
\label{sec:app_algo_rsdrl}

الگوریتم \ref{alg:app_rs_drl_full} شبه‌کد کامل آموزش برخط چارچوب پایه‌ای \lr{RS-DRL} را همراه با تزریق تصادفی پاداش خوش‌بینانه و ذخیره‌سازی تمامی انتقال‌ها در بافر بازپخش تجربه نمایش می‌دهد.

\begin{algorithm}[htbp]
\caption{الگوریتم آموزش چارچوب پایه‌ای \lr{RS-DRL} با شکل‌دهی تصادفی پاداش}
\label{alg:app_rs_drl_full}
\begin{latin}
\small
\begin{algorithmic}[1]
\REQUIRE State space $\mathcal{S}$, action space $\mathcal{A}$, learning rate $\eta$, discount $\gamma$, shaping factor $\rho$, target sync interval $C$
\STATE Initialize main network $Q(s, a; \theta)$ and target network $Q'(s, a; \theta')$ with $\theta' \leftarrow \theta$
\STATE Initialize SPSC-backed replay buffer $\mathcal{D}$ with capacity $N$
\FOR{each episode $= 1, \dots, M$}
    \STATE Observe initial state $s_1$ from managed system
    \FOR{each step $t = 1, \dots, T$}
        \STATE Select $a_t$ via $\epsilon$-greedy policy (random $a_t \in \mathcal{A}$ w.p. $\epsilon$, else $\arg\max_a Q(s_t, a; \theta)$)
        \STATE Execute $a_t$, observe empirical reward $r_{\text{env}}$ and next state $s_{t+1}$
        \STATE Dispatch $(s_t, a_t)$ to UPPAAL-SMC engine and receive $r_{\text{smc}} \in \{0, 1\}$ asynchronously
        \STATE Compute composite reward $r_t \leftarrow r_{\text{env}} \times r_{\text{smc}}$
        \STATE Store $(s_t, a_t, r_t, s_{t+1})$ in $\mathcal{D}$ \COMMENT{Retain both successful and failed transitions}
        \IF{$|\mathcal{D}| \ge B$}
            \STATE Sample random minibatch $\{(s_j, a_j, r_j, s'_{j})\}_{j=1}^B \sim \mathcal{D}$
            \FOR{$j = 1, \dots, B$}
                \IF{$r_j == 0$ \textbf{and} $\text{rand}(0, 1) < \rho$}
                    \STATE $\tilde{r}_j \leftarrow 1.0$ \COMMENT{Optimistic shaping injection for local optima escape}
                \ELSE
                    \STATE $\tilde{r}_j \leftarrow r_j$ \COMMENT{Preserve original continuous or binary reward}
                \ENDIF
                \STATE Compute Bellman target $y_j \leftarrow \tilde{r}_j + \gamma \max_{a'} Q'(s'_j, a'; \theta')$
            \ENDFOR
            \STATE Update $\theta \leftarrow \theta - \eta \nabla_\theta \frac{1}{B} \sum_{j=1}^B (y_j - Q(s_j, a_j; \theta))^2$
        \ENDIF
        \STATE Every $C$ steps synchronize target weights: $\theta' \leftarrow \tau_{\text{soft}}\theta + (1-\tau_{\text{soft}})\theta'$ ($\tau_{\text{soft}}=1.0$)
        \STATE $s_t \leftarrow s_{t+1}$
    \ENDFOR
\ENDFOR
\end{algorithmic}
\end{latin}
\end{algorithm}

\section{الگوریتم سنتز غیربرخط سیاست با هدایت مشخصه‌ها (\texorpdfstring{\sgos}{SG-OS})}
\label{sec:app_algo_sgos}

الگوریتم \ref{alg:app_sgos_full} سازوکار آموزش سیاست بر روی سوابق عملیاتی با عدم توازن بیش از ۹۰٪ داده‌های شکست را مستند می‌سازد.

\begin{algorithm}[htbp]
\caption{الگوریتم سنتز غیربرخط سیاست با هدایت مشخصه‌ها (\sgos)}
\label{alg:app_sgos_full}
\begin{latin}
\small
\begin{algorithmic}[1]
\REQUIRE Offline log $\mathcal{D}_{\text{offline}}$, surrogates $\{M_k\}_{k=0}^{K-1}$, thresholds $\{\tau_k\}_{k=0}^{K-1}$, $\alpha \ge 1$, $\sigma_\Delta$, $\beta$, $\gamma_u$, $\rho_{\text{base}}$, variant $\text{mode} \in \{\text{MU-RS-DRL}, \text{CM-RS-DRL}\}$
\FOR{each training gradient step}
    \STATE Sample minibatch $B \sim \mathcal{D}_{\text{offline}}$
    \FOR{each transition $(s, a, r_{\text{base}}, s', d) \in B$}
        \STATE Compute $\mathcal{K}_{\text{viol}} = \{k \in \{0, \dots, K-1\} \mid M_k(s, a) < \tau_k\}$ and $n_{\text{sat}} = K - |\mathcal{K}_{\text{viol}}|$
        \STATE Compute $\Delta_{\min}(s, a) = \min_k (M_k(s, a) - \tau_k)$ and fractional deficit $\bar{\Delta}_{\text{def}}(s, a)$ via Eq.~(\ref{eq:delta_def_formula})
        \STATE Compute recoverability margin: $\mathcal{M}(s, a) = \left(\frac{n_{\text{sat}}}{K}\right)^\alpha \exp\left(-\frac{\bar{\Delta}_{\text{def}}(s, a)}{\sigma_\Delta}\right)$
        \STATE Compute utility incentive: $\mathcal{U}_{\text{inc}}(a) = \left(\frac{u(a)}{u^{\max}}\right)^\beta$
        \STATE Compute guided probability: $P_{\text{reshape}}(s, a) = \rho_{\text{base}} \cdot \mathcal{M}(s, a) \cdot \left[(1 - \gamma_u) + \gamma_u \mathcal{U}_{\text{inc}}(a)\right]$
        \IF{$\text{mode} == \text{CM-RS-DRL}$}
            \STATE Compute confidence buffer: $\bar{\mathcal{B}}(s, a) = \frac{1}{K} \sum_{k=0}^{K-1} \frac{\max(0, M_k(s, a) - \tau_k)}{1 - \tau_k}$
        \ENDIF
        \IF{$\Delta_{\min}(s, a) \ge 0$}
            \STATE $\tilde{r} \leftarrow 1.0$ if $\text{mode} == \text{MU-RS-DRL}$ else $1.0 + \lambda_{\text{margin}}\bar{\mathcal{B}}(s, a) + \lambda_{\text{util}}\mathcal{U}_{\text{inc}}(a)$
        \ELSIF{$\text{rand}(0, 1) < P_{\text{reshape}}(s, a)$}
            \STATE $\tilde{r} \leftarrow r_{\text{base}}^{\text{MU}} + \mu \mathcal{U}_{\text{inc}}(a)$ if $\text{mode} == \text{MU-RS-DRL}$ else $r_{\text{base}}^{\text{CM}} + \lambda_{\text{margin}}\bar{\mathcal{B}}(s, a) + \lambda_{\text{util}}\mathcal{U}_{\text{inc}}(a)$
        \ELSE
            \STATE $\tilde{r} \leftarrow r_{\text{base}}$ ($= 0$)
        \ENDIF
        \STATE Compute TD target $y = \tilde{r} + (1 - d)\gamma \max_{a' \in \mathcal{A}} Q'(s', a'; \theta')$
    \ENDFOR
    \STATE Update $\theta$ via TD-loss $\mathcal{L}(\theta)$ in Eq.~(\ref{eq:dqn_loss_rsdrlplus})
\ENDFOR
\RETURN Converged policy network $Q^*(s, a; \theta^*)$
\end{algorithmic}
\end{latin}
\end{algorithm}

\section{الگوریتم اعتبارسنجی کشسان با بودجه زمانی مقید (\texorpdfstring{\caepv}{CA-EPV})}
\label{sec:app_algo_caepv}

الگوریتم \ref{alg:app_caepv_full} جزئیات پیاده‌سازی سازوکار پالایش محلی هرزمانی (\pgalr) را تشریح می‌کند.

\begin{algorithm}[htbp]
\caption{الگوریتم اعتبارسنجی کشسان مرحله برنامه‌ریزی با بودجه زمانی مقید (\caepv / \pgalr)}
\label{alg:app_caepv_full}
\begin{latin}
\small
\begin{algorithmic}[1]
\REQUIRE Current state $s_t = (cf_t, of_t)$, policy $Q^*$, surrogates $\{M_k\}_{k=0}^{K-1}$, thresholds $\{\tau_k\}$, stress threshold $\epsilon_{\text{threshold}}$, pre-extracted offline Top-$m$ sensitivity subspace $\text{Top-}m(\mathcal{I}_{\text{XAI}})$, time budget $T_{\text{budget}}$
\STATE $t_{\text{start}} \leftarrow \text{CurrentTime}()$
\STATE \textbf{Stage 1 (Fast Inference):} $a_0 \leftarrow \arg\max_{a \in \mathcal{A}} Q^*(s_t, a)$ \COMMENT{Sub-millisecond nominal proposal}
\STATE \textbf{Stage 2 (Margin \& Stress Gate):} Compute $\Delta_{\min}(s_t, a_0) = \min_k (M_k(s_t, a_0) - \tau_k)$ and $\mathcal{E}(of_t) = \sum_j w_j \widetilde{of}_{t,j}$
\IF{$\Delta_{\min}(s_t, a_0) \ge 0$ \textbf{and} $\mathcal{E}(of_t) \le \epsilon_{\text{threshold}}$}
    \RETURN $a_0$ \COMMENT{Accept via Fast Path without refinement latency}
\ENDIF
\STATE \textbf{Stage 3 (PG-ALR 1-Hop Refinement):} Generate $\mathcal{N}(a_0)$ along pre-extracted $\text{Top-}m(\mathcal{I}_{\text{XAI}})$ via Eq.~(\ref{eq:xai_neighborhood}) \COMMENT{$|\mathcal{N}(a_0)| \le 2m + 1$}
\STATE Evaluate vectorized surrogate compliance $\mathcal{N}_{\text{sat}}(s_t, a_0) = \{ a \in \mathcal{N}(a_0) \mid \Delta_{\min}(s_t, a) \ge 0 \}$ within budget $T_{\text{budget}} - \delta_{\text{guard}}$
\IF{$\mathcal{N}_{\text{sat}}(s_t, a_0) \neq \emptyset$}
    \STATE $a^* \leftarrow \arg\max_{a \in \mathcal{N}_{\text{sat}}(s_t, a_0)} \left[ u(a) + \lambda_{\text{tie}} \bar{\Delta}(s_t, a) \right]$ \COMMENT{Maximize operational utility with tie-breaker}
\ELSE
    \STATE $a^* \leftarrow \arg\max_{a \in \mathcal{N}(a_0)} \sum_{k=0}^{K-1} (M_k(s_t, a) - \tau_k)$ \COMMENT{Margin-maximizing safety fallback}
\ENDIF
\RETURN $a^*$ \COMMENT{Verified adaptation guaranteed by Proposition~2 (Prop.~\ref{prop:safety_preservation})}
\end{algorithmic}
\end{latin}
\end{algorithm}

\newpage

% =========================================================================
% پیوست ج: نتایج تجربی گسترده و جداول فراپارامترها
% =========================================================================
\chapter{نتایج تجربی گسترده و جداول فراپارامترها}
\label{chap:appendix_extended_results}

این پیوست فهرست کامل مشخصات فراپارامتری و جداول ارزیابی‌های آماری تکمیلی را در سناریوهای مأموریتی مختلف گزارش می‌دهد.

\section{فهرست کامل فراپارامترهای یادگیری و شبیه‌سازی}
\label{sec:app_hyperparameters}

جدول \ref{tab:app_full_hyperparams} فهرست کامل مقادیر فراپارامترهای مورد استفاده در آموزش شبکه عصبی و شبیه‌سازی بسترهای \lr{DeltaIoT} و \lr{DARTSim} را نشان می‌دهد.

\begin{table}[htbp]
\centering
\small
\caption{فهرست کامل فراپارامترهای آموزش یادگیری تقویتی عمیق و شبیه‌سازی سیستم‌های هدف (هماهنگ با جدول \ref{tab:hyperparameter_catalog_thesis})}
\label{tab:app_full_hyperparams}
\resizebox{\textwidth}{!}{%
\begin{tabular}{|p{4.2cm}|p{4.0cm}|p{5.8cm}|}
\hline
\textbf{فراپارامتر} & \textbf{مقدار تنظیم‌شده} & \textbf{توضیح فنی و کارکرد} \\ \hline
معماری شبکه در \lr{DARTSim} & \lr{MLP $7 \times 64 \times 64 \times 60$ (ReLU)} & مدل اصلی \lr{SEAMS} روی بردار ۷ بعدی و ۶۰ پیکربندی (و $17 \times 128 \times 128 \times 8$ در پیش‌آزمایش غیربرخط ۸گانه) \\ \hline
معماری شبکه در \lr{DeltaIoT} & \lr{MLP $2 \times 128$ (ReLU)} & تقریب‌گر ارزش در فضاهای ۲۱۶ تا ۴۰۹۶ گزینه شبکه حسگر بی‌سیم \\ \hline
نرخ یادگیری \lr{Adam} ($\eta$) & $10^{-3}$ & گام به‌روزرسانی بهینه‌ساز \lr{Adam} ($\beta_1=0.9, \beta_2=0.999$) در آزمایش‌های اصلی (جدول \ref{tab:hyperparameter_catalog_thesis}) \\ \hline
ضریب تنزیل زمانی ($\gamma$) & $0.99$ & وزن‌دهی به پاداش‌های تجمعی آتی در افق مارکوف ($d=1$ در \lr{DARTSim}) \\ \hline
اندازه دسته آموزش ($B$) & $64$ ($256$ در آزمون هم‌روندی) & تعداد نمونه‌های واکشی‌شده در هر به‌روزرسانی گرادیان \\ \hline
ظرفیت بافر بازپخش ($N$) & $10^5$ کل / $30{,}000$ \lr{DARTSim} / $20{,}000$ \lr{DeltaIoT} & بافر $30{,}000$تایی متوازن از ۵۰ هزار گذار ثبت‌شده (با نرخ شکست $90.98\%$) \\ \hline
دوره همگام‌سازی شبکه هدف ($C$) & $500$ (\lr{DARTSim}) / $200$ (\lr{DeltaIoT}) & همگام‌سازی سخت دوره‌ای ($\tau_{\text{soft}}=1.0$ در پولییاک) \\ \hline
پارامترهای شکل‌دهی (\sgos) & $\rho = 0.30$، $\rho_{\text{base}} = 0.65$، $\alpha = 2.0$، $\sigma_\Delta = 0.05$، $\beta = 1.5$، $\gamma_u = 0.40$ & نرخ شکل‌دهی پایه و فراپارامترهای حاشیه بازیابی و مشوق مطلوبیت ($r_{\text{base}}^{\text{MU}}=0.30, \mu=0.40$) \\ \hline
پارامترهای بافر اطمینان (\cmrsdrl) & $\tau = 0.80$، $\lambda_{\text{margin}} = 0.30$، $\lambda_{\text{util}} = 0.30$، $r_{\text{base}}^{\text{CM}} = 0.70$ & نرمال‌سازی حاشیه اضافه و تقویت پاداش ایمن در مرزهای آستانه‌ای \\ \hline
بُعد حساسیت و تنش (\pgalr) & $m = 3$ ($\delta_{\text{hop}} \le 1$)، $\lambda_{\text{tie}} = 0.05$، $\epsilon_{\text{threshold}} = 0.35$ / $0.65$ & $m=3$ بعد برتر \lr{TreeSHAP}، آستانه تنش $0.35$ (\lr{DARTSim}) و $0.65$ (\lr{DeltaIoT})، سقف زمانی \lr{$<3.5\,\text{ms}$} \\ \hline
\end{tabular}%
}
\end{table}

\section{گزارش ماتریس نتایج آزمون‌های فرضیه آماری}
\label{sec:app_statistical_matrices}

جدول \ref{tab:app_wilcoxon_details} مقادیر دقیق آماره‌های آزمون‌های ناپارامتری (رتبه‌دار ویلکاکسون، مک‌نمار، اندازه اثر وارگا-دلانی $\hat{A}_{12}$ و کوهن $d$) را به تفکیک پرسش‌های پژوهش \lr{RQ1} تا \lr{RQ6} (متناظر با جداول \ref{tab:stat_tests} و \ref{tab:complete_p_values} در فصل \ref{chap:evaluation}) گزارش می‌کند.

\begin{table}[htbp]
\centering
\small
\caption{خلاصه جامع آزمون‌های فرضیه آماری و اندازه‌های اثر به تفکیک پرسش‌های پژوهش (هماهنگ با جداول \ref{tab:stat_tests} و \ref{tab:complete_p_values})}
\label{tab:app_wilcoxon_details}
\resizebox{\textwidth}{!}{%
\begin{tabular}{|p{4.0cm}|p{3.4cm}|c|c|c|p{3.2cm}|}
\hline
\textbf{پرسش پژوهش و جفت مقایسه} & \textbf{نوع آزمون و نمونه} & \textbf{آماره / اختلاف} & \textbf{مقدار دقیق $p$} & \textbf{اندازه اثر} & \textbf{وضعیت معناداری ($\alpha_{\text{adj}}=0.005$)} \\ \hline
\lr{RQ1} — مقدار $Q$: \\ \lr{MU-RS-DRL vs. Baseline DQN} & ویلکاکسون / $t$ روی همگرایی \mbox{($N=5$)} & $Q: +0.64$ در برابر $-5.8$ & $0.0397$ & $d = 3.75$ (استثنایی) & معنادار در مهار افت مقادیر $Q$ \\ \hline
\lr{RQ2} — مطلوبیت \lr{DARTSim}: \\ \lr{MU-RS-DRL vs. Foundational / XDA-II} & ویلکاکسون زوجی \mbox{($N=200$)} & $\Delta d_{\text{norm}} = +0.0385$ / $+0.0595$ & $2.37 \times 10^{-10}$ / $4.12 \times 10^{-10}$ & $\hat{A}_{12} = 0.778$ / $0.795$ & رد فرض صفر \mbox{($p \ll \alpha_{\text{adj}}$)} \\ \hline
\lr{RQ3} / \lr{RQ5} — تأخیر برخط و مهار به زیر \lr{$3.5\,\text{ms}$} & ویلکاکسون / نسبت سرعت \mbox{($N=200$)} & $0.485\,\text{ms}$ در برابر $62.61\,\text{ms}$ (\lr{XDA-II}) & \lr{$< 10^{-15}$} & $129\times\text{--}267\times$ سریع‌تر & برآورده‌سازی قطعی مهلت \lr{$<50\,\text{ms}$} \\ \hline
\lr{RQ4} — انتقال به \lr{DeltaIoTv2} ۴۰۹۶تایی: \\ \lr{RS-DRL\textsuperscript{+} vs. Foundational RS-DRL} & آزمون مقایسه‌ای پیوسته \mbox{($N=30$)} & اتلاف بسته: $8.32\%$ در برابر $10.49\%$ & $4.96 \times 10^{-18}$ & $d = 1.96$ (بزرگ) & رد فرض صفر \mbox{($p \ll \alpha_{\text{adj}}$)} \\ \hline
\lr{RQ6} — انطباق جانشین $v_3$: \\ \lr{MU-RS-DRL vs. FITEST / Random} & مک‌نمار زوجی \mbox{($N=200$)} & $b=137, c=0$ / $b=138, c=0$ & $1.21 \times 10^{-31}$ / $7.29 \times 10^{-32}$ & $\hat{A}_{12} = 0.843$ / $0.845$ & برتری قطعی انطباق \mbox{($p \ll \alpha_{\text{adj}}$)} \\ \hline
\lr{RQ6} — بقای فیزیکی بدون سپر $v_3$: \\ \lr{RS-DRL\textsuperscript{+} vs. Shielded} & مک‌نمار دقیق \mbox{($N=200$)} & $b=49, c=38$ \mbox{($\text{OR}=1.29$)} & $0.284$ & $d = 0.24$ (تناسب) & برابری ایمنی با کاهش $267\times$ تأخیر \\ \hline
\end{tabular}%
}
\end{table}

\newpage

% =========================================================================
% پیوست د: مصنوعات پیاده‌سازی و پکیج بازتولید
% =========================================================================
\chapter{مصنوعات پیاده‌سازی و پکیج بازتولید}
\label{chap:appendix_reproducibility}

به منظور پایبندی به اصول علم باز\LTRfootnote{Open Science} و تکرارپذیری مستقل کلیه شواهد تجربی این رساله دکتری، کلیه کدهای منبع، اسکریپت‌های واسط، ردپاهای محیطی و عملیاتی و فایل‌های کانتینر داکر در دسترس محققان قرار گرفته است.

\section{ساختار مخزن نرم‌افزاری و ماژول‌ها}
\label{sec:app_repo_structure}

مخزن گیت پژوهش شامل ساختار سازمان‌یافته زیر است:

\paragraph{ماژول هسته الگوریتمی:}
پیاده‌سازی پایتون برای \lr{RS-DRL}، \mursdrl، \cmrsdrl و سازوکار پالایش محلی هرزمانی (\pgalr) بر پایه کتابخانه‌های \lr{PyTorch} و \lr{Stable-Baselines3}.

\paragraph{ماژول اتصال به شبیه‌ساز \lr{DeltaIoT}:}
کدهای جاوا و سوکت‌های شبکه جهت ارتباط میان فرآیند یادگیری برخط پایتون و شبیه‌ساز شبکه حسگر با قابلیت شبیه‌سازی تداخلات محیطی و بار ترافیکی پیام‌ها.

\paragraph{ماژول شبیه‌ساز \lr{DARTSim}:}
افزونه‌های \lr{C++}، آداپتور غیربرخط \lr{OfflineDARTSimEnv} و توابع محاسبه مقادیر شاپلی درختی (\lr{TreeSHAP}) در شبیه‌سازی مأموریت ۶ درجه آزادی پرواز.

\paragraph{اسکریپت‌های اعتبارسنجی آماری:}
پردازش سوابق خام، آزمون رتبه‌دار ویلکاکسون، رسم نمودارهای مرز پارتو و کادرهای جعبه‌ای توزیع کیفیت.

\section{دستورالعمل استقرار در محیط کانتینری داکر}
\label{sec:app_docker_setup}

جهت مصون ماندن از تعارض‌های ناشی از نسخه‌های سیستم‌عامل و راه‌اندازها، پکیج بازتولید در قالب یک تصویر داکر مجزا استانداردسازی شده است. مراحل فراخوانی و اجرای کانتینر به شرح زیر است:

\begin{enumerate}
    \item \textbf{کلون کردن مخزن کدهای منبع:}
    \begin{LTR}
    \small
    \begin{verbatim}
    git clone https://github.com/alirezakavianifar/RS-DRL-Dart.git
    cd RS-DRL-Dart
    \end{verbatim}
    \end{LTR}

    \item \textbf{ساخت تصویر داکر (\lr{Docker Build}):}
    \begin{LTR}
    \small
    \begin{verbatim}
    docker build -t rsdrl-reproducibility:v1.0 .
    \end{verbatim}
    \end{LTR}

    \item \textbf{اجرای محیط بازتولید با دسترسی به منابع پردازشی:}
    \begin{LTR}
    \small
    \begin{verbatim}
    docker run -it --cpus="16" --memory="32g" \
           -v $(pwd)/results:/workspace/results \
           rsdrl-reproducibility:v1.0
    \end{verbatim}
    \end{LTR}
\end{enumerate}

\section{دستورات بازتولید آزمایش‌ها و آزمون‌های فرضیه}
\label{sec:app_reproduction_commands}

پس از ورود به محیط کانتینر، اجرای سناریوهای ارزیابی با دستورات خط فرمان زیر امکان‌پذیر است:

\subsubsection*{بازتولید نتایج بستر \lr{DeltaIoT} (\lr{RQ1} تا \lr{RQ4})}
\begin{LTR}
\small
\begin{verbatim}
python run_deltaiot_benchmarks.py --instances v1 v1.1 v2 --seeds 30
\end{verbatim}
\end{LTR}
این دستور شبیه‌سازی همگرایی و هرس فضا را روی هر سه مقیاس شبکه حسگر اجرا کرده و سوابق نتایج را در مسیر \lr{\texttt{results/deltaiot/}} ذخیره می‌نماید.

\subsubsection*{بازتولید ۲,۰۰۰ پرواز مأموریتی \lr{DARTSim} (\lr{RQ1}، \lr{RQ5} و \lr{RQ6})}
\begin{LTR}
\small
\begin{verbatim}
python run_dartsim_eval.py --scenarios 200 --flights-per-scenario 1 \
                           --methods RS-DRL-Plus MU-RS-DRL CM-RS-DRL \
                                     Foundational-RS-DRL Baseline-DQN \
                                     Shielded XDA-II-SHAP XDA-PDP \
                                     NSGA-III FITEST
\end{verbatim}
\end{LTR}
این دستور $2{,}000$ پرواز حلقه‌بسته با دینامیک ۶ درجه آزادی (۱۰ روش $\times$ ۲۰۰ سناریو؛ ۱ پرواز به ازای هر جفت روش-سناریو) را شبیه‌سازی کرده و خروجی‌های بقا، تأخیر میلی‌ثانیه‌ای و مصرف حافظه را در فایل \lr{\texttt{results/}}\allowbreak\lr{\texttt{dartsim\_metrics.csv}} ثبت می‌نماید.

\subsubsection*{تولید نمودارها و ماتریس آزمون‌های آماری ویلکاکسون}
\begin{LTR}
\small
\begin{verbatim}
python generate_statistical_report.py --input-dir results/
\end{verbatim}
\end{LTR}
این اسکریپت آزمون‌های آماری را اجرا کرده، ضرایب پیوسته و اندازه‌های اثر کوهن را محاسبه می‌نماید و نمودارهای مرز پارتو و جعبه‌ای را با فرمت \lr{PDF} با وضوح بالا بازتولید می‌کند.

\newpage



\ifstandalone
\bibliographystyle{ieeetr-fa}
\bibliography{References,MyReferences}
\fi
\end{document}
